\setlength{\footskip}{8mm}

\chapter{Conclusion and Recommendations}
\label{ch:conclusion}

\section{Contributions}
\begin{enumerate}
    \item Using hydrological, meteorological, and temporal variables to train and evaluate water level forecasting models; benchmarking numerous machine learning models and conventional baselines.
    \item Water-level prediction using an end-to-end, fully deep learning LSTM pipeline that deals with the requirement for complicated hydrodynamic models, like those found in nationwide forecasting systems.
    \item A feature ablation analysis that measures the importance of each feature group supports a comprehensive feature set including lags, rolling statistics, difference features, and cycling encoding.
    \item Reliable real-world behavior is ensured by an effective time-aware evaluation and advanced error analysis to measure performance under various temporal circumstances.
    \item Risk-level categorization system that improves interpretability by mapping predictions into flood risk classifications based on bank and bed reference levels.
    \item Ready for production: a useful tool for monitoring and prediction with live weather data, dynamic diagrams, and automated 24-hour forecasts.
    \item \textbf{Weather Impact Analysis:} Demonstrated that incorporating meteorological variables (rainfall, pressure, humidity) improves prediction accuracy by \textbf{24\%}, serving as critical leading indicators for water level changes.
\end{enumerate}

\section{Recommendations and Future Work}

Future work will focus on several key areas to enhance the system's capabilities:
\begin{itemize}
    \item \textbf{Models:} Explore Transformer architectures and ensemble methods to further improve accuracy.
    \item \textbf{Data:} Integrate upstream station data and real-time weather forecasts instead of historical data.
    \item \textbf{Horizons:} Develop multi-horizon forecasting capabilities (6h, 12h, 24h, 48h, 72h).
    \item \textbf{Deployment:} Implement cloud deployment with automated retraining pipelines.
    \item \textbf{Uncertainty:} Incorporate probabilistic predictions with confidence intervals to better inform decision-making.
\end{itemize}
